\subsection{Ablation studies}
\label{sec:analysis}

\paragraph{Data augmentation and correction are complementary.}
Augmentation improves fixed-parameter ResNet success on PushObj by 17.4 points and adaptive success by 5.4 points; ECP then adds a further 11.1 and 2.9 points on the same checkpoints. The gains stack because augmentation improves the experience available during training, while ECP improves the objective used during execution. The augmentation benefit is not specific to our backbone: under matched training, it raises DINO-WM with CEM from 24.6\% to 48.9\%, improving all seven shapes (\cref{tab:dinowm-aug}).



\begin{table}[t]
\centering
\caption{\textbf{Component ablation on PushObj} ($\gamma_{\rm hi}=0.15$). Success (\%) on the same 350 tasks and 30-replan budget. $S$ denotes memory-based stagnation regularization; $P$ denotes temporary proprioceptive reweighting. Both checkpoints are offline-pretrained ResNet world models. Frozen fixes their parameters; AdaJEPA updates them online. Native omits ECP. Row maxima are bold.}
\label{tab:components}
\small
\setlength{\tabcolsep}{3.2pt}
\begin{tabular}{llrrrr}
\toprule
Checkpoint & Updates & Native & $S$ & $P$ & $S{+}P$\\
\midrule
Released & Frozen &43.7&44.9&\textbf{60.0}&59.7\\
Released & AdaJEPA &69.4&71.4&72.9&\textbf{73.1}\\
Data aug. & Frozen &61.1&62.6&\textbf{72.3}&\textbf{72.3}\\
Data aug. & AdaJEPA &74.9&76.0&\textbf{77.1}&76.9\\
\bottomrule
\end{tabular}
\end{table}

\paragraph{Both components improve native planning.}
Across both checkpoints and both update rules, stagnation regularization and temporary goal reweighting each improve native PushObj success (\cref{tab:components}). Reweighting supplies the largest manipulation gains: fixed-parameter ResNet success rises from 43.7\% to 60.0\% with the released model and from 61.1\% to 72.3\% with the augmented-data model. Combined correction improves the four native baselines by 16.0, 3.7, 11.1, and 2.0 points, respectively, and matches reweighting alone to within one task in every setting. The components serve distinct recovery needs: goal-component emphasis drives the PushObj gains, while spatial stagnation memory redirects navigation in Maze.

\begin{table}[t]
\centering
\caption{\textbf{Trigger-threshold ablation on PushObj.} Augmented-data ResNet checkpoint with AdaJEPA online updates and $S{+}P$, 350 tasks, 30 replans, $\gamma_{\rm lo}=\gamma_{\rm hi}/3$. Triggered counts episodes with at least one detected stall; First is the median replan of the first trigger. Rescue/harm counts native failures recovered and native successes lost.}
\label{tab:threshold}
\small
\setlength{\tabcolsep}{4pt}
\begin{tabular}{lrrrr}
\toprule
$\gamma_{\rm hi}$ & Success (\%) & Triggered & First & Rescue / harm\\
\midrule
Native &74.9&---&---&---\\
0.15 &76.9&90&14&8 / 1\\
0.30 &77.4&101&12&11 / 2\\
0.45 &\textbf{77.7}&108&11&11 / 1\\
0.60 &77.4&111&10&11 / 2\\
\bottomrule
\end{tabular}
\end{table}

\paragraph{ECP is robust to the trigger threshold.}
\Cref{tab:threshold} sweeps the upper threshold over a fourfold range. A larger threshold detects stalls in more episodes (90 to 111) and earlier (median first trigger from replan 14 to 10), yet every setting improves on native planning by 7 to 10 tasks with at most two regressions. Success peaks at $\gamma_{\rm hi}=0.45$, which we adopt for the main PushObj results. With ResNet parameters fixed, $\gamma_{\rm hi}=0.15$ and $0.45$ both solve 253 of 350 tasks, each with 39 rescues and no regression. ECP therefore does not depend on a finely tuned threshold.

\begin{table}[t]
\centering
\caption{\textbf{Triggered correction versus permanent reweighting and additional search} ($\gamma_{\rm hi}=0.15$). PushObj planning with the offline-pretrained, augmented-data ResNet world model held fixed. Permanent weighting uses the same $10\times$ multiplier as ECP. Rescue/harm counts are paired changes relative to native MPC.}
\label{tab:timing}
\small
\begin{tabular}{lrr}
\toprule
Planning variant & Success (\%) & Rescue / harm\\
\midrule
Native MPC &61.1&---\\
Permanent weighting &66.3&50 / 32\\
Lookahead search &63.1&34 / 27\\
ECP + triggered lookahead &69.4&29 / 0\\
ECP (triggered $S{+}P$) &\textbf{72.3}&39 / 0\\
\bottomrule
\end{tabular}
\end{table}

\paragraph{Conditional intervention preserves successful behavior.}
With the same $10\times$ multiplier, triggered correction solves 253 tasks versus 232 for permanent weighting (\cref{tab:timing}). It retains all 214 native successes and recovers 39 failures, while permanent weighting loses 32 native successes. Conditioning the objective on a detected stall therefore preserves successful behavior.

\paragraph{Changing the objective outperforms searching longer.}
Lookahead search reaches 63.1\% with 27 regressions; combining lookahead with ECP reaches 69.4\%. ECP alone reaches 72.3\% without regressions under the same execution budget. Additional action search does not resolve a poorly directed objective.

\paragraph{Joint latent memory supports scene-aware correction.}
With AdaJEPA and reweighting disabled, joint visual/proprioceptive memory reaches 75.4--76.0\% across three strengths, versus at most 75.1\% for proprioceptive-only memory (\cref{tab:memory-strength}). Visual state distinguishes configurations that share an agent position. All variants preserve the native successes.

\begin{table}[t]
\centering
\caption{\textbf{Stagnation memory on Maze.} Frozen holds the offline-pretrained ResNet world model fixed; AdaJEPA updates it online. Native omits ECP; Ours adds memory-only $S$. Success requires position error $<0.5$ on 100 tasks with a 30-replan budget.}
\label{tab:maze}
\small
\begin{tabular*}{0.94\columnwidth}{@{\extracolsep{\fill}}llr@{}}
\toprule
Model updates & Planner & Success (\%) $\uparrow$\\
\midrule
Frozen & Native &48.0\\
& Ours &\textbf{50.0}\\
AdaJEPA & Native &63.0\\
& Ours &\textbf{66.0}\\
\bottomrule
\end{tabular*}
\end{table}

\begin{table}[t]
\centering
\caption{\textbf{Rope planning with DINO-WM.} Both rows keep the pretrained model fixed and predict/execute three pushes per plan. DINO-WM (ours) adds dynamic $S$ after five warm-up pushes. Success means Chamfer distance (CD) $<0.3$, over 20 targets with at most 20 pushes each.}
\label{tab:rope}
\small
\begin{tabular*}{0.94\columnwidth}{@{\extracolsep{\fill}}lrr@{}}
\toprule
Method & CD $\downarrow$ & Success (\%) $\uparrow$\\
\midrule
DINO-WM (native) &0.755&50.0\\
DINO-WM (ours) &\textbf{0.515}&\textbf{60.0}\\
\bottomrule
\end{tabular*}
\end{table}

\paragraph{Correction complements continued adaptation.}
Gated memory retains all 63 native Maze successes while solving three additional tasks. Immediate correction also reaches 66 successes but incurs four regressions; fixed delay reaches only 61. The loss gate coordinates correction with model learning (\cref{sec:maze-details}).

\subsection{Qualitative analysis of recovery}
\label{sec:qualitative}

In the PushObj example (\cref{fig:teaser}), the fixed-parameter ResNet controller repeatedly moves near the goal while leaving the pusher misplaced. ECP reduces pusher error from 45.3 to 10.2 pixels and succeeds at replan 17, after sharing the first 15 plans. \Cref{fig:execution-recovery} adds shifted PushT and Rope examples.

In Maze, AdaJEPA travels 45.1 units while repeatedly traversing the same corridor. Gated memory redirects the trajectory and reaches the goal at replan 27; the runs share their first six plans. With the ResNet model held fixed, memory reaches the goal at replan 28 after sharing the first 13 plans with the failed baseline. These cases show goal-component completion and escape from revisited regions. All solid trajectories follow replayed physics steps that match the logged replan endpoints.
